%-----------------------------------LICENSE------------------------------------%
%   This file is part of Mathematics-and-Physics.                              %
%                                                                              %
%   Mathematics-and-Physics is free software: you can redistribute it and/or   %
%   modify it under the terms of the GNU General Public License as             %
%   published by the Free Software Foundation, either version 3 of the         %
%   License, or (at your option) any later version.                            %
%                                                                              %
%   Mathematics-and-Physics is distributed in the hope that it will be useful, %
%   but WITHOUT ANY WARRANTY; without even the implied warranty of             %
%   MERCHANTABILITY or FITNESS FOR A PARTICULAR PURPOSE.  See the              %
%   GNU General Public License for more details.                               %
%                                                                              %
%   You should have received a copy of the GNU General Public License along    %
%   with Mathematics-and-Physics.  If not, see <https://www.gnu.org/licenses/>.%
%------------------------------------------------------------------------------%
%       Author: Ryan Maguire                                                   %
%       Date:   2026/02/23                                                     %
%------------------------------------------------------------------------------%
\documentclass{article}
\usepackage{graphicx}
\setlength{\parindent}{0em}
\setlength{\parskip}{0em}
\title{On the Need for a New Sampling Requirement in RSS Reconstructions}
\author{Ryan Maguire}
\date{\today}
\graphicspath{{../images/}}
\begin{document}
    \maketitle
\    The one-dimensional Fresnel transform may be obtained from the
    Rayleigh-Huygens equations by applying the stationary phase method.
    The normalized complex diffraction profile $\hat{T}$ may be written
    in terms of the complex transmittance $T$ via:

    \begin{equation}
        \hat{T}(\rho_{0})
        =
        \frac{1-i}{2F}
        \int_{0}^{\infty}
            T(\rho)
            \exp\left(
                i\psi(\rho,\,\phi_{s},\,\rho_{0},\,\phi_{0})
            \right)\,\textrm{d}\rho
    \end{equation}

    where $\psi$ is the Fresnel kernel, and $\phi_{s}$ is the staionary
    azimuth angle, the solution to $\partial\psi/\partial\phi=0$, and $F$
    is the Fresnel scale.
    When the opening angle is large, the classic quadratic Fresnel
    approximation may be used to simplify this calculation:

    \begin{equation}
        \hat{T}(\rho_{0})
            \approx
            \frac{1-i}{2F}
            \int_{0}^{\infty}
                T(\rho)
                \exp\left(
                    \frac{i\pi}{2}\left(
                        \frac{\rho-\rho_{0}}{F}
                    \right)^{2}
                \right)\,\textrm{d}\rho
    \end{equation}

    This is in the form of a convolution, and an inverse formula may be
    obtained using the convolution theorem:

    \begin{equation}
        T(\rho)
        \approx
        \frac{1+i}{2F}
        \int_{0}^{\infty}
            \hat{T}(\rho_{0})
            \exp\left(
                -\frac{i\pi}{2}\left(
                    \frac{\rho-\rho_{0}}{F}
                \right)^{2}
            \right)\,\textrm{d}\rho_{0}
    \end{equation}

    In practice we restrict our attention to a finite window $W$ centered
    about $\rho$ so that the integral becomes:

    \begin{equation}
        T(\rho)
        \approx
        \frac{1+i}{2F}
        \int_{\rho-W/2}^{\rho+W/2}
            w(\rho-\rho_{0})
            \hat{T}(\rho_{0})
            \exp\left(
                -\frac{i\pi}{2}\left(
                    \frac{\rho-\rho_{0}}{F}
                \right)^{2}
            \right)\,\textrm{d}\rho_{0}
    \end{equation}

    where $w$ is the tapering function. Using a simple $\delta$ function
    centered at $\rho_{*}$ for $T$, the diffraction profile $\hat{T}$ becomes:

    \begin{equation}
        \begin{array}{rcl}
            \displaystyle
            \hat{T}(\rho_{0})
            \!\!
            &\approx&
            \!\!
            \displaystyle
            \frac{1-i}{2F}
            \int_{0}^{\infty}
                T(\rho)
                \exp\left(
                    \frac{i\pi}{2}\left(
                        \frac{\rho-\rho_{0}}{F}
                    \right)^{2}
                \right)\,\textrm{d}\rho\\[1.5em]
            \!\!
            &=&
            \!\!
            \displaystyle
            \frac{1-i}{2F}
            \int_{0}^{\infty}
                \delta(\rho-\rho_{*})
                \exp\left(
                    \frac{i\pi}{2}\left(
                        \frac{\rho-\rho_{0}}{F}
                    \right)^{2}
                \right)\,\textrm{d}\rho\\[1.5em]
            \!\!
            &=&
            \!\!
            \displaystyle
            \frac{1-i}{2F}
            \exp\left(
                \frac{i\pi}{2}\left(
                    \frac{\rho_{*}-\rho_{0}}{F}
                \right)^{2}
            \right)
        \end{array}
    \end{equation}

    Feeding this into the inverse transform for $T$ produces:

    \begin{equation}
        \begin{array}{rcl}
            \displaystyle
            T(\rho)
            \!\!
            &\approx&
            \!\!
            \displaystyle
            \frac{1+i}{2F}
            \int_{\rho-W/2}^{\rho+W/2}
                \hat{T}(\rho_{0})
                \exp\left(
                    -\frac{i\pi}{2}\left(
                        \frac{\rho-\rho_{0}}{F}
                    \right)^{2}
                \right)\,\textrm{d}\rho_{0}\\[1.5em]
            \!\!
            &=&
            \!\!
            \displaystyle
            \frac{1+i}{2F}
            \int_{\rho-W/2}^{\rho+W/2}
                \frac{1-i}{2F}
                \exp\left(
                    \frac{i\pi}{2}\left(
                        \left(
                            \frac{\rho_{*}-\rho_{0}}{F}
                        \right)^{2}
                        -
                        \left(
                            \frac{\rho-\rho_{0}}{F}
                        \right)^{2}
                    \right)
                \right)\,\textrm{d}\rho_{0}\\[1.5em]
            \!\!
            &=&
            \!\!
            \displaystyle
            \frac{%
                \exp\left(%
                    \frac{i\pi}{2F^{2}}\left(\rho_{*}^{2}-\rho^{2}\right)%
                \right)
            }{2F^{2}}
            \int_{\rho-W/2}^{\rho+W/2}
                \exp\left(
                    \frac{i\pi}{F^{2}}\rho_{0}\left(\rho-\rho_{*}\right)
                \right)\,\textrm{d}\rho_{0}\\[1.5em]
            \!\!
            &=&
            \!\!
            \displaystyle
            \frac{W}{2F^{2}}\,
            \textrm{sinc}\left(
                \frac{W}{2F^{2}}(\rho-\rho_{*})
            \right)
            \exp\left(
                \frac{i\pi}{2}\left(
                    \frac{\rho-\rho_{*}}{F}
                \right)^{2}
            \right)
        \end{array}
    \end{equation}

    This follows directly from MTR86. The input is a delta response, and the
    recovered profile is a sinc function. The resolution may then be defined
    as half the width of the main lobe:

    \begin{equation}
        \Delta{R}
        =
        \frac{2F^{2}}{W}
    \end{equation}

    That is, if a resolution $\Delta$ has been requested, then the window
    width of integration must be at least:

    \begin{equation}
        W
        =
        \frac{2F^{2}}{\Delta{R}}
    \end{equation}

    From this we may obtain an estimate for how finely the data must be
    sampled for us to adequately perform the numerical integration needed for
    the inverse transform.
    \par\hfill\par
    The diffraction profile $\hat{T}$ from a $\delta$ response has highly
    oscillatory terms. A simple bound for the sampling rate is that
    two samples per wave are recorded (Shannon's sampling theorem).
    For $\exp(\frac{i\pi{x}^{2}}{2F^{2}})$,
    the zeros in the imaginary part occur at $F\sqrt{2n}$, where $n$ is an
    integer. The wavelength of the $n^{\footnotesize\textrm{th}}$ wave is hence:

    \begin{equation}
        \lambda_{n}
        =
        F\left(
            \sqrt{2n+4}-\sqrt{2n}
        \right)
    \end{equation}

    Since $x$ is bounded by $W/2$, solving $F\sqrt{2n}=W/2$ shows that
    $n$ is bounded by $W^{2}/8F^{2}$. The smallest wave in the window has
    a wavelength of:

    \begin{equation}
        \lambda_{min}
        =
        F\left(
            \sqrt{4+\frac{W^{2}}{4F^{2}}}
            -
            \frac{W}{2F}
        \right)
        =
        \frac{-W+\sqrt{16F^{2}+W^{2}}}{2}
    \end{equation}

    The maximum allowed distance between samples should be half of this.

    \begin{equation}
        \Delta{x}
        =
        \frac{-W+\sqrt{16F^{2}+W^{2}}}{4}
    \end{equation}

    Note, this has an asymptotic expansion (Laurent series) in $W$ given by:

    \begin{equation}
        \Delta{x}
        =
        \frac{2F^{2}}{W}
        -
        \frac{8F^{4}}{W^3}
        +
        \frac{64F^{6}}{W^{5}}
        -
        \cdots
    \end{equation}

    As an alternative, we simply find where the derivative of the phase is
    maximum. For $\phi(x)=\pi{x}^{2}/2F^{2}$, this is achieved at the endpoints
    where the slope is $\omega_{\textrm{max}}=\pi{W}/2F^{2}$.
    The smallest period is hence $2\pi/\omega_{\textrm{max}}$, and we need
    two samples per period, giving us:

    \begin{equation}
        \Delta{x}
        =
        \frac{2F^{2}}{{W}}
    \end{equation}

    This is the leading term for the asymptotic series found using the
    previous definition of $\lambda_{\textrm{min}}$, meaning for large windows
    (high resolutions) the two methods agree quite well.
\end{document}
